\section{Introduction}
\label{sec:introduction}

Modern cities increasingly deploy sensor-based platforms to monitor environmental
conditions (air quality, temperature, noise levels) and feed real-time data to operators
who must react to degraded conditions. As these platforms adopt Service-Oriented
Architectures (SOA) built on microservices, they inherit the security challenges of distributed
systems: an expanded attack surface, heterogeneous trust boundaries, and complex
inter-service communication that must be secured without introducing prohibitive latency.

This project presents the design, implementation, and evaluation of a threat detection
and prevention system for a Smart City Environmental Monitoring Platform built as a SOA.

\subsection{Scope}

The system addresses four threat classes, with two sub-variants of the spam threat:

\begin{enumerate}[leftmargin=*]
  \item \textbf{HTTP Flood (DoS):} volumetric single-source flood targeting the
        operator portal, aiming to blind city personnel before or during a physical event.
  \item \textbf{Advanced Persistent Threat (APT):} stealthy, multi-stage intrusion
        combining reconnaissance and targeted exploitation to gain unauthorised access
        to internal platform resources.
  \item \textbf{Alert Flood:} abuse of legitimate platform functionality to generate
        a high volume of automated notifications, degrading operator situational awareness.
  \item \textbf{Drug Content Spam:} unsolicited email containing illicit pharmaceutical
        content, targeting the platform's mail infrastructure as a delivery channel.
  \item \textbf{Botnet DDoS:} distributed volumetric attack orchestrated across multiple
        compromised nodes, overwhelming the target service through coordinated flooding.
\end{enumerate}

The APT and Alert Flood threats form a multi-stage attack chain: initial intrusion
enables subsequent service abuse, mirroring the lateral movement and impact phases of
real-world adversarial campaigns.

The remainder of this report is organised as follows. Section 2 reviews the relevant
state of the art. Section 3 describes the service architecture. Section 4 presents the
risk management analysis and mitigation strategies. Sections 5 through 9 detail the
detection and response approach for each threat class. Section 10 evaluates the overall
solution, its thresholds, and its scalability. Section 11 draws conclusions and outlines
future work.
